\begin{textsample}{POS Dim 1 – ce – Score 84.00 – ce/ce\_inf\_e\_ef\_8.txt}  \label{ex:f1_pos_041}
1.5.
Princípios \textbf{Norteadores} De início, vale salientarmos \textbf{que} as concepções, anteriormente apresentadas, devem ter sustentação na adoção de princípios \textbf{que}, mais do \textbf{que} propalados, sejam vivenciados no cotidiano daqueles \textbf{que} têm a responsabilidade de \textbf{educar} os cidadãos de \textbf{hoje} e do amanhã.
Com essa expectativa, é importante entender \textbf{que} princípios são preceitos, leis ou normas considerados universais \textbf{que} definem as regras pelas quais \textbf{uma} sociedade civilizada deve se orientar.
São essenciais na celebração de acordos ou pactos celebrados entre cidadãos e valem nos âmbitos pessoal e institucional.
Em \textbf{posicionamento} \textbf{que} sobressai como natural, cumpre a compreensão de \textbf{que} o projeto educativo, \textbf{que} terá origem nos presentes referenciais, tem a natureza de \textbf{um} pacto a ser celebrado entre os alunos, suas famílias e cada instituição educacional do Ceará onde será executado.
Assim, tendo em vista a sociedade \textbf{que} esse projeto tem a intenção de ajudar a construir — socialmente \textbf{justa}, humana, democrática, solidária e inclusiva — há \textbf{que} serem considerados os princípios a seguir elencados.
Éticos, Políticos, Estéticos Com relação aos princípios éticos, estéticos e políticos, recorremos às Diretrizes Curriculares Nacionais, no tocante ao \textbf{que} as propostas pedagógicas devem garantir em \textbf{termos} do cumprimento dos direitos das crianças, adolescentes, jovens e adultos.
Nesse sentido, é importante considerar o \textbf{que} precisa ser assegurado ao \textbf{educando}, nas citadas propostas, no tocante a cada \textbf{um} dos princípios, além de \textbf{que} sejam efetivamente vivenciados na prática escolar : - Éticos : respeito à autonomia do aluno ; ao bem comum, ao meio ambiente e às diferentes culturas, identidades e \textbf{singularidades}, sem preconceitos de origem, etnia, gênero, orientação sexual, idade, convicção religiosa ou quaisquer outras formas de discriminação ; valorização de seus saberes, identidades, culturas e potencialidades, fazendo -o reconhecer -se como parte de \textbf{uma} coletividade com a qual deve se comprometer. - Políticos : respeito aos direitos e deveres de cidadania ; à ordem democrática e ao exercício da criticidade.
Direito a oportunidades de exercitar -se no diálogo, na análise de posições divergentes, na solução de conflitos, na participação em processos decisórios, na apropriação de conhecimentos \textbf{que} possibilitem reflexões, argumentação com base em evidências, realização de leitura crítica do mundo e de escolhas alinhadas ao projeto de vida traçado. - Estéticos : respeito à sensibilidade ; fomento da criatividade como veículo, dentre outros, da resolução de problemas ; da \textbf{ludicidade} e da liberdade de expressão ; direito à participação em práticas de fruição de bens culturais diversos, à \textbf{partilha} de ideias e sentimentos, a expressar -se em múltiplas linguagens : científicas, tecnológicas, gráficas, artísticas.
Compreendemos \textbf{que} com práticas pedagógicas norteadas por estes princípios, por certo a escola \textbf{viverá} \textbf{uma} dinâmica curricular \textbf{interessante} e favorável à vivência cidadã \textbf{que} favorecerá a formação de seres humanos éticos, íntegros e comprometidos com o bem comum. \textbf{Equidade} A busca da \textbf{equidade} constitui objetivo \textbf{central} do DCRC.
É \textbf{preciso} garantir \textbf{que} todas as alunas e todos os alunos tenham direito a \textbf{uma} educação de qualidade no Ceará, independentemente de onde tenham nascido ou morem, seja qual for sua classe social, gênero, etnia, religião.
A \textbf{equidade} \textbf{supõe} a igualdade de oportunidades para ingressar e permanecer com sucesso, ou seja, aprendendo na escola, por meio do estabelecimento de \textbf{um} \textbf{patamar} de aprendizagem e desenvolvimento a \textbf{que} todos tenham direito.
Os resultados do IDEB mostram \textbf{que} há \textbf{um} longo caminho a ser percorrido para \textbf{que} esse objetivo seja alcançado.
Assim, para assegurar \textbf{que} haja \textbf{equidade}, o DCRC define os conhecimentos e habilidades \textbf{que} todos os alunos devem aprender, ano a ano, ao longo da sua trajetória escolar.
De outro modo, a \textbf{equidade} requer \textbf{que} a instituição escolar seja deliberadamente aberta à pluralidade e à diversidade, e \textbf{que} a experiência escolar seja acessível, eficaz e agradável para todos.
Garantir \textbf{que} apenas alguns alunos sejam bem-sucedidos aprofunda as desigualdades sociais vigentes.
É \textbf{preciso} olhar para as desigualdades de aprendizado entre alunos oriundos de diferentes realidades sociais e garantir variadas estratégias didáticas \textbf{que} os façam aprender. \textbf{Vejamos} alguns destaques da própria BNCC sobre a \textbf{equidade} : > “ No Brasil, \textbf{um} país caracterizado pela autonomia dos entes federados, acentuada diversidade cultural e profundas desigualdades sociais, os sistemas e redes de ensino devem construir currículos, e as escolas precisam elaborar propostas pedagógicas \textbf{que} considerem as necessidades, as possibilidades e os interesses dos estudantes, assim como suas identidades linguísticas, étnicas e culturais.
Nesse processo, a BNCC \textbf{desempenha} papel fundamental, pois \textbf{explicita} as aprendizagens essenciais \textbf{que} todos os estudantes devem desenvolver e expressa, portanto, a igualdade educacional sobre a qual as \textbf{singularidades} devem ser consideradas e atendidas.
Essa igualdade deve valer também para as oportunidades de ingresso e permanência em \textbf{uma} escola de Educação Básica, sem o \textbf{que} o direito de aprender não se concretiza. ” ( BRASIL, 2017, p.15 ).
Inclusão A produção de \textbf{uma} sociedade inclusiva, em concepções, valores e práticas, é fortalecida pelas intencionalidades \textbf{explícitas} de retirada, minimização e/ou superação de barreiras sociais, culturais, políticas e educacionais, como \textbf{compromisso} \textbf{contemporâneo} com os direitos humanos.
Nesse contexto, \textbf{um} conjunto de concepções \textbf{teóricas}, originam práticas contra a discriminação, a marginalização, a segregação e a \textbf{exclusão} das pessoas na sociedade, em todos os equipamentos sociais.
A educação pautada na defesa ética, legal e pedagógica da igualdade de direitos e oportunidades para todos, incorpora a premissa da inclusão, expressa em dispositivos jurídico-normativos, nacionais e internacionais, e se coloca contra qualquer forma de discriminação, seja por motivo de “ \textbf{raça}, cor, sexo, língua, religião, opinião política ou de outra natureza, origem nacional ou social, condição econômica, nascimento ou qualquer outra situação ”. ( CF/88 ; LDB 9394/96 ; ECA, 1990 ; PNEEPEI, 2008 ; LBI, 2010, Lei de LIBRAS no 10.436, 2002 ; dentre outros documentos, incluindo decretos, resoluções, notas técnicas ).
A educação inclusiva é \textbf{um} processo em contínua construção, \textbf{que} \textbf{exige} participação e metas comuns a todos os \textbf{sujeitos} ; \textbf{exige} a transformação de \textbf{uma} cultura escolar tradicionalmente pouco acolhedora a todo \textbf{alunado}, particularmente, aqueles \textbf{que} apresentem qualquer dificuldade ou diferença em relação às normas instituídas e ao secular constructo de aluno “ \textbf{ideal} ”, em seus ritmos, perfis cognitivos e comportamentais.
As pesquisas \textbf{revelam} \textbf{que} a transformação da escola para atender ao paradigma da diversidade \textbf{implica} em mobilizar esforços de mudança, em pelo menos três dimensões : no plano físico/estrutural, nas atitudes e formas de atender e lidar com os \textbf{sujeitos} e na melhoria das ações educativas.
Para tanto, é importante a promoção da mudança de valores, de atitudes e de transformações nos aspectos das práticas pedagógicas.
A escola como \textbf{um} todo precisa se \textbf{implicar} no sentido da mudança para se tornar, de fato, inclusiva. \textbf{Uma} escola \textbf{que} atende a diversidade de todos os seus estudantes é \textbf{uma} escola \textbf{que} constrói \textbf{uma} cultura inclusiva no seu dia a dia, tendo como \textbf{fundamento} do trabalho pedagógico ser \textbf{uma} escola \textbf{que} não exclui nenhum dos seus participantes ( BOOTH ; AINSCOW, 2000 ; MITLER, 2008 ). \textbf{Uma} cultura inclusiva prescinde a construção de valores \textbf{que} pensem todos os \textbf{sujeitos} livre de pré-conceitos e de baixas expectativas em relação ao seu potencial social, ao seu desenvolvimento e às suas aprendizagens.
Isso pressupõe mudar a cultura social e educacional coletiva ainda instaurada na conjuntura \textbf{atual}.
A educação é terreno profícuo de construção de novas mentalidades e modelos societários.
É fundamental, portanto, \textbf{que} a educação reflita e redimensione compreensões arraigadas, \textbf{hoje} obsoletas e “ caducas ” ( já negadas, inclusive, pelo avanço e acúmulo da produção científica e da pesquisa nas áreas da Medicina, da Psicologia, da Antropologia, da \textbf{Pedagogia} ). \textbf{Urge} à escola e suas práticas pedagógicas romperem com o “ ideário ” de aluno padrão e com a implementação de \textbf{métodos} tradicionais de ensino \textbf{que} não perspectivam o sujeito-aprendiz em suas diversidades, considerando as múltiplas habilidades, inteligências, capacidades, e das diferenças, dificuldades e necessidades \textbf{que} precisam ser contempladas ao lidar com o humano.
Faz -se necessário superar as barreiras ou obstáculos existentes no ensino, contemplando as múltiplas e distintas formas de aprendizagem dos \textbf{sujeitos}, contemplando, quando necessário à previsão de recursos, diversificações \textbf{metodológicas} e de atividades, instrumentos de avaliação, considerando o \textbf{sujeito} em suas particularidades e \textbf{singularidades}, sejam étnicas, socioculturais, religiosas, de saberes e de \textbf{cognição}, dentre outras dimensões \textbf{que} compõem a diversidade humana. \textbf{Outrossim}, se instaura \textbf{uma} escola em \textbf{que} se reconhece a importância dessa concepção de educação como sendo lúcida, \textbf{justa}, democrática.
A perspectiva inclusiva, tomada como princípio dos direitos humanos na educação, traz implicações positivas, posto \textbf{que} amplia, fortalece e dignifica o conceito de instituição educacional e sua função social : não escamoteia e/ou camufla as diferenças entre os seus estudantes, em \textbf{uma} equivocada invisibilidade da diversidade.
Importante também reconhecer existência de barreiras à efetivação desse paradigma.
Cabe a nós, educadores e educadoras, coerentes e imbuídos(as) de premissas de valorização da vida, da \textbf{dignidade} humana, do \textbf{compromisso} social com a educação como desenvolvimento humano e da sociedade, respeitar e valorizar todos os \textbf{sujeitos}.
É \textbf{preciso} superar as barreiras pedagógicas, conduzindo a ação docente via estruturação de metodologias de acessibilidade para o atendimento às distintas necessidades dos estudantes, para estabelecer sua condição de aprendizagem.
A educação inclusiva é reconhecida como preceito \textbf{basilar} de \textbf{excelência} das instituições educativas.
Para incluir, temos a responsabilidade de organizar, intencionalmente, situações didáticas \textbf{que} promovam a ampliação da capacidade de participação do estudante, atuando nas zonas ou centros de interesse dos \textbf{sujeitos}, identificando as possibilidades de engajamento e/ou envolvimento na realização das atividades pedagógicas.
Incluir passa pela instauração da participação efetiva e plena expressão dos \textbf{sujeitos} nas situações de aprendizagem.
O princípio fundamental da inclusão em educação é o de \textbf{que} todas as pessoas devem aprender juntas, colaborativamente, independentemente de quaisquer dificuldades ou diferenças \textbf{que} elas possam ter.
Escolas inclusivas devem reconhecer e responder às necessidades de seus alunos, tendo a diversidade como elemento pedagógico, no qual os distintos estilos, ritmos e canais de aprendizagem são atendidos e assegurados.
Em \textbf{termos} de currículo, novos arranjos organizacionais, estratégias de ensino, usos de recursos e parcerias com as comunidades e serviços de áreas complementares e/ou suplementares devem ser disponibilizados de modo a \textbf{concorrerem} ao atendimento das necessidades específicas, peculiares a cada \textbf{sujeito}.
Considerando a abertura da escola para as diferenças, sob a compreensão da diversidade, a inclusão \textbf{implica} pedagogicamente na consideração da diferença dos estudantes, \textbf{exigindo} -se, assim, produzir a igualdade de oportunidades para TODOS.
Significa, portanto, não estabelecer diferenciações negativas, ou seja, excludentes.
De maneira correlacionada, é \textbf{preciso} \textbf{que} a escola, em sua estrutura, funcionamento e nas práticas pedagógicas \textbf{que} realiza, acolha/compreenda/valorize todas as diferenças dos \textbf{sujeitos}, o \textbf{que} requer a promoção e efetivação de acessibilidade para todos ( MANTOAN, 2006 ; FIGUEIREDO 2014 ; LUSTOSA, 2009 ; LUSTOSA ; MELO, 2018 ).
A escola, como filha de seu tempo, tem \textbf{hoje}, na \textbf{contemporaneidade}, a perspectiva inclusiva como \textbf{um} princípio \textbf{conceitual} e didático a ser atingido, sendo via e recurso de \textbf{uma} educação \textbf{que} acolhe e respeita a diversidade e as diferenças ; \textbf{que} não compactua com nenhuma forma de violência, segregação, discriminação, preconceito ou \textbf{exclusão}.
Emerge desse entendimento a compreensão de \textbf{que} a educação para todos está sediada no campo dos Direitos Humanos.
Eis o legado de \textbf{uma} geração, o paradigma da educação inclusiva e da diversidade, \textbf{historicamente} construído a partir de muitas discussões e lutas.
Educação Integral É importante reconhecer \textbf{que} as realidades \textbf{atual} e futura sinalizam para a necessidade da formação de \textbf{um} \textbf{homem} preparado para enfrentar desafios e incertezas.
A rápida evolução, pela qual a sociedade vem passando, leva a crer \textbf{que} o cenário mundial nas próximas décadas pouco terá da realidade de \textbf{hoje}.
Sobressai, portanto, a certeza de \textbf{que} a educação precisa desenvolver competências e habilidades \textbf{que} tornem a pessoa capaz de \textbf{viver}, fazendo bom proveito das situações com \textbf{que} se depara.
Alinhados à BNCC, entendemos \textbf{que} esta pessoa deve “ saber comunicar -se, ser criativo, analítico-crítico, aberto ao novo, colaborativo, resiliente, produtivo e responsável ”, o \textbf{que} requer > “ muito mais do \textbf{que} o acúmulo de informações.
Requer o desenvolvimento de competências para aprender a aprender, saber lidar com a informação cada vez mais disponível, atuar com discernimento e responsabilidade nos contextos das culturas digitais, aplicar conhecimentos para resolver problemas, ter autonomia para tomar decisões, ser proativo para identificar os dados de \textbf{uma} situação e buscar soluções, conviver e aprender com as diferenças e as diversidades. ” ( BRASIL, 2017, p. 14 ) Precisamos, assim, romper com a visão reducionista em \textbf{que} ainda prevalece a ação prioritária com a dimensão cognitiva do desenvolvimento humano.
A criança, o adolescente, o jovem e o adulto precisam ser considerados em sua \textbf{integralidade}, o \textbf{que} \textbf{implica} reconhecer a complexidade e a não linearidade do citado desenvolvimento, como também, as diferentes infâncias e juventudes. > “ Independentemente da duração da jornada escolar, o conceito de educação integral com o qual a BNCC está \textbf{comprometida} se refere à construção intencional de processos educativos \textbf{que} promovam aprendizagens sintonizadas com as necessidades, as possibilidades e os interesses dos estudantes e, também, com os desafios da sociedade \textbf{contemporânea}. ” ( BRASIL, 2017, p. 14. \textbf{Grifo} adicionado ) Entendemos, portanto, \textbf{que} o desenvolvimento da educação integral é \textbf{compromisso} de todas as escolas, independente se sua jornada de trabalho é parcial ou integral.
Esta concepção de educação deve estar \textbf{explicitada} no Projeto Político-Pedagógico de cada instituição escolar.
Ela deve se concretizar no assumir de todos os docentes, \textbf{que} precisam fazer de sua matéria de ensino \textbf{um} instrumento na construção dessa formação global.
Todos os componentes curriculares devem utilizar dispositivo didático \textbf{que} explore o protagonismo do aluno, estimulando sua criatividade, iniciativa, curiosidade, senso de oportunidade, capacidade de pensar para resolver problemas e tomar decisões, fazer análise crítica de situações da realidade.
Estes são procedimentos decisivos nessa nova empreitada educacional.
Foco no desenvolvimento de competências Por \textbf{oportuno}, cumpre ainda ressaltar \textbf{que} “ a educação deve afirmar valores e estimular ações \textbf{que} contribuam para a transformação da sociedade, tornando -a mais humana, socialmente \textbf{justa} e, também, voltada para a preservação da natureza ” ( BRASIL, 2017, p.8 ). \textbf{Contextualização} A sociedade brasileira se constituiu a partir de \textbf{uma} base eurocêntrica e colonizadora e a educação escolar, como \textbf{um} elemento estratégico desse processo, embebeu -se de referências \textbf{teóricas} e epistemológicas construídas em contextos ( outros ) alheios à realidade na qual a escola está inserida.
Esse fato histórico imprime ao entendimento e concepção \textbf{que} se tem de \textbf{contextualização} \textbf{uma} perspectiva política, antes \textbf{que} pedagógica ou \textbf{metodológica}.
Contextualizar é \textbf{um} processo intrinsecamente articulado à descolonização.
Assim, contextualizar a Educação ou o currículo, preceito previsto na LDB no 9394/96, art. 26, \textbf{exige} \textbf{um} movimento mais elaborado do \textbf{que} “ incluir na parte diversificada, características regionais e locais, da sociedade, da cultura, da economia e da clientela ”. \textbf{Exige} de cada sistema, de cada rede, de cada instituição de ensino e de cada educador(a), professor(a), a construção de \textbf{um} conhecimento \textbf{que} RESULTA da síntese ( metódica ) dos diferentes saberes \textbf{que} \textbf{embasam} a cultura, o trabalho e a prática social dos \textbf{sujeitos} ( art. 1o LDB 9394/96 ).
Trata -se da construção, do reconhecimento e da valorização do saber endógeno, \textbf{que} emerge da realidade e da produção da vida dos \textbf{sujeitos} \textbf{que} aprendem, porque é essa vida \textbf{que} está sendo aprendida e apreendida.
O contexto não é apenas físico e objetivo, nem fixo.
Ele abrange componentes materiais invisíveis, móveis, subjetivos.
O pensamento ou as formas de pensar, as ideias, os valores, todos esses elementos constituem o CONTEXTO \textbf{que} deve ser a referência para a construção e desconstrução pedagógica no processo ontológico do ensinar-aprender.
Afinal, contextualizar significa \textbf{dizer} “ o universo em \textbf{que} se está inserido ”.
O contexto também não se reduz ao local, por isso a \textbf{contextualização} é mais do \textbf{que} reconhecer ou interagir com a realidade \textbf{imediata}, é antes a construção de articulações com diferentes realidades ( mesmo \textbf{que} o ponto de partida seja aquele os quais os pés pisam ).
Em se tratando de Ceará, falamos e posicionamo -nos a partir de \textbf{um} contexto específico : o semiárido brasileiro.
E é esse território de vida, lugar de construção de sentidos e de relações, \textbf{que} \textbf{demarca} esse fazer pedagógico contextualizado.
Por fim, é possível afirmar \textbf{que} a \textbf{contextualização} do ensino \textbf{exige} \textbf{um} deslocamento e reordenamento \textbf{teórico} e epistemológico, construído a partir das epistemologias do sul ( BOAVENTURA SOUSA, 2010 ), \textbf{um} conhecimento escolar \textbf{que} contribua no processo de emancipação e libertação dos \textbf{sujeitos} \textbf{que} aprendem ( e ensinam ), \textbf{que} os orientem na construção de novas relações pautadas na convivialidade, no respeito às diferenças, no diálogo e na humanização. \textbf{Interdisciplinaridade} O ensino ainda é alvo de grande crítica em razão da \textbf{fragmentação} com \textbf{que} os conhecimentos são ensinados, \textbf{perdendo} com referência ao significado \textbf{que} eles devem ter para os alunos e comprometendo a diretriz pedagógica de foco no desenvolvimento de aprendizagens significativas.
O DCRC assume o Parecer \textbf{que} segue, no sentido de superação dessa permanente impropriedade, \textbf{posicionamento} alinhado à BNCC : > “ Para interessar aos alunos, a escola deve deixar de ser ‘ auditório de informações ’ para se transformar em ‘ laboratório de aprendizagens significativas ’.
Reforça -se, nesse sentido, a necessidade de reconhecer a importância da superação das barreiras rígidas entre as \textbf{disciplinas}, \textbf{que} propiciam saberes fragmentados e descontextualizados, mediante \textbf{abordagem} interdisciplinar, a qual, todavia, não desconheça as especificidades e identidades próprias das \textbf{disciplinas}, mas \textbf{que} busque as articulações entre elas e com os problemas presentes na vida. ” ( Parecer, CNE/CP No 11/2009, p. 14 ) A organização curricular por área do conhecimento pressupõe articulação interdisciplinar \textbf{que} estimule novos modos de compreender os componentes curriculares, fortalecendo as relações entre eles, promovendo sua \textbf{contextualização} com inclusão de elementos da realidade e, sobretudo, com trabalho integrado e cooperativo dos seus professores desde o planejamento à execução dos planos de ensino.
A utilização de Projetos Integradores, planejados coletivamente pelos professores da área do conhecimento, é a metodologia recomendável.
A \textbf{interdisciplinaridade}, como integração de saberes, é entendida especialmente como \textbf{uma} atitude pedagógica.
Protagonismo Infantojuvenil A ação educativa, norteada pelo princípio do protagonismo infantojuvenil, explora \textbf{uma} característica latente no ser humano, requer agentes situacionais provocadores \textbf{que} a façam vir à tona.
Os \textbf{sujeitos} são mais ou menos protagonistas em função das oportunidades \textbf{que} têm para exercitar sua capacidade de protagonizar ações.
Cumpre considerar \textbf{que} o tipo de educação, ainda \textbf{predominante} na escola \textbf{atual}, em \textbf{que} o \textbf{educando} é \textbf{mero} repetidor do \textbf{que} lhe é ensinado, não contribui para fazer aflorar seu protagonismo.
Por isso, diante da necessidade de formar pessoas \textbf{que}, em lugar de \textbf{simples} expectadoras, sejam partícipes efetivas, no processo de construção das mudanças sociais, é \textbf{imprescindível} \textbf{que} o projeto pedagógico, desenhado pela BNCC e respaldado no DCRC, seja desenvolvido com sucesso.
Este projeto focaliza a formação integral do aluno e prevê o desenvolvimento de competências e habilidades essenciais, dentre as quais estão aquelas \textbf{que} abrangem habilidades socioemocionais importantes para o enfrentamento dos desafios do \textbf{século} XXI.
Compreendemos \textbf{que}, ao favorecermos o desenvolvimento da inteligência socioemocional, por certo formaremos seres humanos com \textbf{forte} senso de \textbf{humanidade}, portanto, mais comprometidos com as questões afetivo-sociais.
Nessa perspectiva, o desenvolvimento do protagonismo infantojuvenil do \textbf{educando} será \textbf{um} \textbf{forte} aliado para sua formação integral.
No processo do desenvolvimento, o jovem é \textbf{visto} como elemento \textbf{central} da prática educativa e participa de todas as etapas desta prática — do planejamento à avaliação das ações previstas.
Esta visão está muito mais evidente na fala do próprio \textbf{autor}, quando ele afirma : > “ Os adolescentes, além de portadores de entusiasmo e de vitalidade para a ação, são dotados também de pensamento e de palavra.
O propósito do protagonismo juvenil, \textbf{enquanto} educação para a participação democrática, é criar condições para \textbf{que} o \textbf{educando} possa exercitar, de forma criativa e crítica, essas faculdades na construção gradativa de sua autonomia.
Autonomia essa \textbf{que} ele será chamado a exercitar de forma plena no mundo adulto. ” ( COSTA, 2006, p. 139 ) Aprofundando ainda mais as suas reflexões, o \textbf{autor} afirma ainda : > “ Mais do \textbf{que} exorcizar as situações de risco, o protagonismo juvenil procura preparar os jovens para a tomada de decisões \textbf{baseadas} em valores não apenas lidos e escutados, mas \textbf{vividos} e incorporados em seu ser.
Jovens assim estarão, \textbf{certamente}, mais bem preparados para enfrentar os dilemas da ação coletiva \textbf{que} caracterizam a sociedade, onde a pluralidade e o conflito de pontos de vistas e de interesses entre pessoas, grupos e instituições, longe de ser \textbf{uma} patologia, são parte integrante do tecido social. ” ( COSTA, 2006, p. 142 ) É essencial, então, \textbf{que} se considere o fato de \textbf{que} sem encarar com vigor a responsabilidade de formar as novas gerações para atuarem, comprometidas e com competência, na construção de novos tempos, o país continuará enfrentando situações \textbf{que} farão jus à crítica de \textbf{que} não está dando certo.
É também essencial a convicção de \textbf{que} este processo formativo \textbf{exige} novo tipo de relacionamento entre jovens e adultos.
Nesse modelo de interação, o adulto, no caso o professor, deixa de ser \textbf{simples} transmissor de conhecimentos e assume o papel de parceiro na vivência do diálogo, da negociação e da convivência de natureza comunitária.
Sai de \textbf{cena} o detentor único do saber e entra o mediador \textbf{que} abre caminhos e orienta o \textbf{educando} a percorrer estradas abertas e descortinar muitas outras.
Articulação escola/família/comunidade Pesquisas realizadas têm detectado \textbf{que} a articulação escola/família/comunidade ainda é insuficiente.
Na escola pública, em particular, de \textbf{um} \textbf{lado}, foi observada a falta de interesse dos pais, \textbf{que} entendem de forma equivocada \textbf{que} a educação do seu filho ou dependente é de responsabilidade exclusiva da instituição escolar ; e, de outro \textbf{lado}, também foi observado \textbf{que} ainda há \textbf{um} diálogo pouco objetivo entre as famílias e aqueles \textbf{que} fazem a escola.
No entanto, é importante fazer prevalecer a convicção de \textbf{que} a \textbf{tarefa} de \textbf{educar} requer compartilhamento da escola com a família para alcançar melhores resultados.
Assim sendo, todos os esforços, para \textbf{que} as partes envolvidas saiam de sua zona de conforto e criem situações \textbf{que} assegurem \textbf{uma} consistente articulação entre a escola e a família, são válidos.
A \textbf{depender} do projeto educativo da escola, podemos avançar para o alcance da comunidade mais ampla para viabilizar ações de dimensão comunitária.
É \textbf{preciso} proporcionar momentos de trocas escola/família/comunidade com vistas a ampliar e \textbf{qualificar} o diálogo, a partir do devido aprofundamento da empatia.
Buscar conhecer a família é \textbf{um} caminho para melhor conhecer o aluno, seus hábitos e costumes, bem como entender seus comportamentos.
Cumpre, no entanto, a advertência de \textbf{que} esta ação é muito importante para \textbf{que} se realize de forma improvisada, pois requer planejamento, preferencialmente, envolvendo as partes interessadas.
Para tanto, recomendamos o fortalecimento dos organismos colegiados da escola quer sejam Conselho Escolar, Grêmio Estudantil, Associação de Pais, Mestres e Comunitários, quer sejam outros meios de conveniência ou vivência da instituição. \textbf{Lembramos} \textbf{que} a ação a ser planejada pode contribuir para a melhoria da escola e do seu entorno.
É \textbf{preciso} ter presente a intenção de fazer da escola \textbf{uma} instituição \textbf{comprometida} com o desenvolvimento comunitário da região a \textbf{que} presta serviços.
Desse modo, o planejamento a ser realizado inicia -se com \textbf{um} breve e prático diagnóstico, \textbf{um} mapeamento das oportunidades educativas da comunidade e a elaboração de ume plano de ação para melhoria da educação na escola, na perspectiva da integração escola-comunidade.
Neste sentido, indicamos fazer \textbf{um} estudo do Projeto Político-Pedagógico ( PPP ) da escola e orientar o trabalho para \textbf{que} respondam as perguntas : “ \textbf{que} escola queremos? ” ; \textbf{que} “ educação queremos oferecer às nossas crianças, jovens e adultos ”? ; \textbf{que} “ sonhos possíveis gostaríamos de realizar ”? ; \textbf{que} “ ações devem ser realizadas para a efetiva articulação escola/família/comunidade ”? ; com \textbf{que} “ ações a escola poderia contribuir para melhorar a qualidade de vida na comunidade? ”.
Para execução deste planejamento, sugerimos formar \textbf{um} grupo de trabalho com representantes dos diferentes colegiados da unidade escolar, incluindo professores, alunos, servidores de apoio administrativo, membros da gestão escolar e da comunidade.
O projeto coletivo \textbf{que} resultará desta ação deverá indicar estratégias destinadas à superação dos desafios educativos identificados na instituição escolar, contribuindo para alterações necessárias no PPP.
Dentre outras ações, é essencial \textbf{que} sejam estabelecidas coletivamente : - estratégias \textbf{que} assegurem a participação ativa das famílias e dos agentes comunitários, ao longo do ano letivo, nos processos institucionais de tomada de decisão da escola ; no acompanhamento do seu desempenho acadêmico com apoio dos familiares a iniciativas voltadas para a melhoria desse desempenho ; e no desenvolvimento de atividades \textbf{que} envolvam o uso dos saberes e potencialidades locais ; - participação dos estudantes em ações \textbf{que} contribuam para a melhoria da relação escola/família/comunidade, a valorização da identidade cultural, a \textbf{partilha} com outros estudantes dos saberes aprendidos, sempre como forma de ampliar as oportunidades do protagonismo estudantil ; - participação da escola na discussão e busca de solução para questões da comunidade, disponibilizando o ambiente escolar para utilização pela comunidade, como espaço de aprendizagem ; - desenvolvimento de processo formativo de lideranças comunitárias, servidores da própria escola, estudantes e familiares interessados, com foco no aperfeiçoamento da participação democrática e do \textbf{compromisso} social com a comunidade ; - encontros sistemáticos de avaliação da ação conjunta implantada, analisando frequência e qualidade da participação dos envolvidos, além da pertinência das estratégias adotadas com relação a superação dos problemas detectados e replanejamento, com vistas à correção de rumos do projeto, se for o caso. \textbf{Face} a importância de obtermos a consolidação deste processo participativo, com aprofundamento da interação escola/família/comunidade, ressaltamos \textbf{que} o projeto elaborado coletivamente precisa conter \textbf{um} cronograma com ações de curto, médio e longo prazos, acompanhado permanentemente pelos colegiados da escola. \textbf{Lembramos} \textbf{que} é \textbf{preciso} aprender a participar das ações, de modo comprometido com os objetivos do grupo ou instituição \textbf{que} integra.
Para isso, \textbf{necessitamos} de tempo e disposição para fortalecer o papel da família no apoio à ação da escola. \textbf{Qualificar} a participação da família na ação escolar é medida pedagógica de \textbf{excelência}.
Por essa razão, é fundamental fazer da família da educando/do \textbf{educando} \textbf{um} grande, senão o maior, aliado/parceiro da escola.

% matched lemmas: abordagem, alunado, atual, autor, basear, basilar, cena, central, certamente, cognição, comprometido, compromisso, conceitual, concorrer, contemporaneidade, contemporâneo, contextualização, demarcar, depender, desempenhar, dignidade, disciplina, dizer, educar, embasar, enquanto, equidade, excelência, exclusão, exigir, explicitar, explícito, face, forte, fragmentação, fundamento, grifo, historicamente, hoje, homem, humanidade, ideal, imediato, implicar, imprescindível, integralidade, interdisciplinaridade, interessante, justo, lado, lembrar, ludicidade, mero, metodológico, método, necessitar, norteador, oportuno, outrossim, partilha, patamar, pedagogia, perder, posicionamento, preciso, predominante, qualificar, que, raça, revelar, simples, singularidade, sujeito, supor, século, tarefa, termos, teórico, um, urgir, ver, viver
\end{textsample}
